\section{Future Work}

Several limitations point to concrete next steps. The clearest is data coverage: front-right appears in only three positive recordings and front-both is rarer still (Table~\ref{tab:dataset}), which shows up as the weakest per-leg attribution ($\mathrm{F1}=0.636$ for FR in Table~\ref{tab:perleg}) and as the dominant source of LORO variance ($\mathrm{F1}=0.807\pm0.142$, Fig.~\ref{fig:loro}); collecting more front-right and front-both interactions should close that attribution gap and tighten the spread across folds. Every detection number we report is offline replay of recorded logs, so an on-robot field study that measures live detection quality and end-to-end latency on the Go2's onboard computer, rather than the dev-CPU benchmark of Fig.~\ref{fig:latency}, is the natural validation. Because no force-instrumented ground truth exists, our severity output is a calibrated index rather than a measured quantity; if instrumented ground truth becomes available, the intensity head could be trained under direct supervision instead of a physics pseudo-target. Two smaller directions round this out: the ablation in Fig.~\ref{fig:ablation} shows the inertial channels are nearly redundant (removing them leaves LORO F1 unchanged to two decimals), so a leaner channel set could trim per-window cost for even tighter latency margins; and since temperature scaling de-saturated the probabilities and slightly improved the Brier score but worsened ECE (Fig.~\ref{fig:reliability}), a better-calibrated confidence signal remains an open goal.
